\section{矩阵的初等变换}

\subsection{初等变换的概念}

\begin{definition}
下面三种变换称为矩阵的\textbf{初等行变换}：
\begin{enumerate}
\item 对换矩阵的两行（对换$i,j$两行，记作$r_i\leftrightarrow r_j$）；
\item 以数$k\neq0$乘矩阵的某一行（第$i$行乘$k$，记作$r_i\times k$）；
\item 把矩阵的某一行乘数$k$加到另一行上去（第$j$行的$k$倍加到第$i$行，记作$r_i+kr_j$）。
\end{enumerate}
把上述定义中的“行”换成“列”，即得\textbf{初等列变换}的定义（记作$c_i\leftrightarrow c_j$、$c_i\times k$、$c_i+kc_j$）。初等行变换与初等列变换统称\textbf{初等变换}。
\end{definition}

三种初等变换都是可逆的，且其逆变换是同一类型的初等变换：
\[
r_i\leftrightarrow r_j\ \text{的逆变换是其本身};\qquad
r_i\times k\ \text{的逆变换是}\ r_i\times\frac{1}{k};\qquad
r_i+kr_j\ \text{的逆变换是}\ r_i+(-k)r_j.
\]

\begin{definition}
如果矩阵$A$经有限次初等变换变成矩阵$B$，就称矩阵$A$与$B$\textbf{等价}，记作$A\cong B$。
\end{definition}

等价是矩阵之间的一种关系，具有三条性质：反身性（$A\cong A$）、对称性（$A\cong B\Rightarrow B\cong A$）、传递性（$A\cong B,B\cong C\Rightarrow A\cong C$）。等价关系把全体$m\times n$矩阵分成若干等价类，同一等价类的矩阵“形状不同而本质相同”（秩相同，见§2）。

\subsection{行阶梯形矩阵与行最简形矩阵}

\begin{definition}
称矩阵$A$为\textbf{行阶梯形矩阵}，如果它满足：
\begin{enumerate}
\item 若有零行（元素全为零的行），则零行全部位于非零行的下方；
\item 各非零行左起第一个非零元的列标，由上至下严格递增。
\end{enumerate}
\end{definition}

行阶梯形矩阵的“台阶”逐行右移、逐行下降，形如阶梯。若再要求每个非零行的首非零元都是$1$、且这些$1$是它们所在列的唯一非零元，则称为\textbf{行最简形矩阵}。

\begin{theorem}
任何矩阵$A$总可以只经初等行变换化为行阶梯形矩阵，并可进一步化为行最简形矩阵。
\end{theorem}

\textbf{化行阶梯形的步骤}（自上而下，逐列推进）：
\begin{enumerate}
\item 在剩下的子矩阵中选一个非零元作为主元（一般选绝对值小、易计算者，或先用$r_i\leftrightarrow r_j$把非零元换到当前行）；
\item 用$r_j+kr_i$把主元\textbf{下方}同列的元素全部化为$0$（若主元不为$1$，可先用$r_i\times k$化为$1$以避免分数）；
\item 对下方的子矩阵重复上述过程，直至剩下的行全为零行。
\end{enumerate}

\textbf{化行最简形}：在行阶梯形的基础上，\textbf{自下而上}用$r_i+kr_j$把每个首非零元\textbf{上方}的元素也化为$0$，再用$r_i\times k$把首非零元化为$1$。

\corepoint{
（1）初等\textbf{行}变换与初等\textbf{列}变换一般不能混用：解线性方程组只允许行变换（列变换会改变未知量的对应关系）；求秩、求标准形时行、列变换都可用。
（2）化行阶梯形时，主元所在列下方必须全为$0$，否则不是阶梯形；主元位置选择不当会使计算量剧增。
（3）行最简形的首非零元是$1$且其所在列“上下都干净”，这是读出方程组的解、求逆矩阵的关键形态。
}

\subsection{初等矩阵}

\begin{definition}
由单位矩阵$E$经过\textbf{一次}初等变换得到的方阵，称为\textbf{初等矩阵}。初等矩阵共有三类：
\begin{enumerate}
\item 把单位矩阵的第$i,j$两行（列）对换，记作$E(i,j)$；
\item 以数$k\neq0$乘单位矩阵的第$i$行（列），记作$E(i(k))$；
\item 以数$k$乘单位矩阵的第$j$行加到第$i$行（或以$k$乘第$i$列加到第$j$列），记作$E(i,j(k))$。
\end{enumerate}
\end{definition}

\begin{theorem}
设$A$是一个$m\times n$矩阵，则
\begin{enumerate}
\item 对$A$施行一次初等\textbf{行}变换，相当于在$A$的\textbf{左边}乘一个相应的$m$阶初等矩阵；
\item 对$A$施行一次初等\textbf{列}变换，相当于在$A$的\textbf{右边}乘一个相应的$n$阶初等矩阵。
\end{enumerate}
\end{theorem}

初等矩阵的性质：
\begin{enumerate}
\item 初等矩阵都是可逆矩阵，且逆矩阵仍是同类型的初等矩阵：
\[
E(i,j)^{-1}=E(i,j),\qquad E(i(k))^{-1}=E\left(i\left(\tfrac{1}{k}\right)\right),\qquad E(i,j(k))^{-1}=E(i,j(-k));
\]
\item 初等矩阵的转置仍是初等矩阵：$E(i,j)'=E(i,j)$，$E(i(k))'=E(i(k))$，$E(i,j(k))'=E(j,i(k))$；
\item $|E(i,j)|=-1$，$|E(i(k))|=k$，$|E(i,j(k))|=1$，故初等矩阵的行列式都不为$0$。
\end{enumerate}

\begin{theorem}
方阵$A$可逆的充要条件是：存在有限个初等矩阵$P_1,P_2,\cdots,P_k$，使
\[
A=P_1P_2\cdots P_k.
\]
\end{theorem}

由此得\textbf{初等变换法求逆矩阵}：设$A$为$n$阶可逆矩阵，由$P_1P_2\cdots P_sA=E$得$P_1P_2\cdots P_sE=A^{-1}$。两式表明：若一系列初等行变换把$A$化成$E$，则同样的行变换把$E$化成$A^{-1}$。于是对分块矩阵$(A\mid E)$施行初等行变换，当$A$化成$E$时，右边的$E$就化成了$A^{-1}$：

\keyformula{(A\mid E)\xrightarrow{\ \text{初等行变换}\ }(E\mid A^{-1})}

类似地，若$A$可逆，则
\[
(A\mid B)\xrightarrow{\text{行变换}}(E\mid A^{-1}B),\qquad
\begin{pmatrix}A\\B\end{pmatrix}\xrightarrow{\text{列变换}}\begin{pmatrix}E\\BA^{-1}\end{pmatrix}.
\]
这两条一次性解出了矩阵方程$AX=B$与$XA=B$，不必先求$A^{-1}$。

\corepoint{
$(A\mid E)\to(E\mid A^{-1})$过程中\textbf{只能做行变换}，一旦用了列变换结果必错。若中途出现全零行，说明$A$不可逆（$|A|=0$），应立即停止并判定“不可逆”。
}

\section{矩阵的秩}

\subsection{秩的概念}

在$m\times n$矩阵$A$中，任取$k$行与$k$列（$k\le\min\{m,n\}$），位于这些行列交叉处的$k^2$个元素按原位置构成的$k$阶行列式，称为$A$的一个\textbf{$k$阶子式}。$A$共有$C_m^kC_n^k$个$k$阶子式。

\begin{definition}
矩阵$A$中非零子式的最高阶数称为矩阵$A$的\textbf{秩}，记作$R(A)$（或$r(A)$、$\mathrm{rank}(A)$）。
\end{definition}

规定零矩阵的秩为$0$。由定义直接得：
\begin{enumerate}
\item $0\le R(A)\le\min\{m,n\}$；
\item $R(A)=0\iff A=O$；
\item $R(A)\ge r\iff A$中至少有一个$r$阶子式不为零；
\item $R(A)<r\iff A$中所有$r$阶子式全为零；
\item $R(A)=r\iff A$有$r$阶非零子式，而所有$r+1$阶子式（若存在）全为零。
\end{enumerate}

\begin{theorem}
矩阵经初等变换后，其秩不变。
\end{theorem}

证明思路：只需证一次初等变换不改变秩。对“倍乘”“倍加”变换，非零子式的对应关系可直接搬运；对“对换”变换，子式至多变号。反过来，因初等变换可逆，秩也不会增大，故秩不变。

\begin{corollary}
$R(A)=r$的充要条件是$A$经初等变换可化为
\[
\begin{pmatrix}E_r&O\\O&O\end{pmatrix}_{m\times n},
\]
称之为$A$的\textbf{等价标准形}。即存在$m$阶可逆矩阵$P$与$n$阶可逆矩阵$Q$，使$PAQ=\begin{pmatrix}E_r&O\\O&O\end{pmatrix}$。
\end{corollary}

\subsection{求矩阵的秩}

由定义求秩要检查大量子式，实际计算用初等变换：

\begin{enumerate}
\item 用初等行变换（也可兼用列变换）把$A$化为\textbf{行阶梯形矩阵}；
\item 行阶梯形矩阵中\textbf{非零行的行数}就是$R(A)$。
\end{enumerate}

道理是：行阶梯形矩阵的非零行数等于其非零子式的最高阶数，而初等变换不改变秩。

\corepoint{
（1）$R(A)$是一个\textbf{数}（非负整数），不是矩阵；写“$R(A)=0$”表示零矩阵，不是“$A=0$”之外的意思。
（2）$R(A)=n$（$n$为阶数）$\iff|A|\neq0\iff A$可逆$\iff A\cong E_n$，这四条等价说法必须能互相转换。
（3）求秩时若用了列变换，最后读“非零行数”仍然正确（秩不变），但若目的是解方程组，列变换会打乱未知量次序。
}

\subsection{秩的性质与不等式}

设$A$为$m\times n$矩阵，$B$为$n\times p$矩阵，$P,Q$为可逆矩阵，则
\begin{enumerate}
\item $R(A')=R(A)$（转置不改变秩）；
\item $R(PA)=R(AQ)=R(PAQ)=R(A)$（乘可逆矩阵不改变秩）；
\item $R(A+B)\le R(A)+R(B)$（$A,B$同型）；
\item $R(AB)\le\min\{R(A),R(B)\}$；
\item $R(AB)\ge R(A)+R(B)-n$（Sylvester不等式）；
\item 特别地，若$AB=O$，则$R(A)+R(B)\le n$；
\item $\max\{R(A),R(B)\}\le R(A,B)\le R(A)+R(B)$，其中$(A,B)$是把$A,B$按列拼成的分块矩阵；
\item $R\begin{pmatrix}A&O\\O&B\end{pmatrix}=R(A)+R(B)$；
\item 若$A$为$n$阶方阵且$A^2=A$，则$R(A)+R(A-E)=n$。
\end{enumerate}

\keyformula{R(AB)\le\min\{R(A),R(B)\},\qquad R(AB)\ge R(A)+R(B)-n,\qquad AB=O\Rightarrow R(A)+R(B)\le n}

\corepoint{
$R(A+B)\le R(A)+R(B)$一般\textbf{不能取等}，$R(AB)\le\min\{R(A),R(B)\}$同样不能取等；而$AB=O$时$R(A)+R(B)\le n$是“取等”的常见情形（当且仅当$A$的零空间维数恰好等于$R(B)$）。含未知参数的秩讨论，方法是：先化行阶梯形，令某个主元恰为$0$，分别讨论参数取值。
}

\subsection{满秩矩阵}

设$A$为$m\times n$矩阵：
\begin{itemize}
\item 若$R(A)=m$，称$A$为\textbf{行满秩矩阵}；
\item 若$R(A)=n$，称$A$为\textbf{列满秩矩阵}；
\item 若$R(A)=\min\{m,n\}$，简称$A$为\textbf{满秩矩阵}；
\item $n$阶方阵$A$满秩$\iff R(A)=n\iff|A|\neq0\iff A$可逆。
\end{itemize}

满秩矩阵的等价标准形形态：
\[
A\ \text{列满秩}\iff A\cong\begin{pmatrix}E_n\\O\end{pmatrix};\qquad
A\ \text{行满秩}\iff A\cong\begin{pmatrix}E_m&O\end{pmatrix}.
\]

\section{线性方程组的解}

\subsection{解的判定定理}

设有$m$个方程、$n$个未知量的线性方程组
\[
\begin{cases}
a_{11}x_1+a_{12}x_2+\cdots+a_{1n}x_n=b_1,\\
a_{21}x_1+a_{22}x_2+\cdots+a_{2n}x_n=b_2,\\
\quad\cdots\cdots\cdots\cdots\\
a_{m1}x_1+a_{m2}x_2+\cdots+a_{mn}x_n=b_m,
\end{cases}
\]
记系数矩阵$A=(a_{ij})_{m\times n}$，未知量列向量$X=(x_1,\cdots,x_n)'$，常数项列向量$\beta=(b_1,\cdots,b_m)'$，则方程组可写成矩阵形式$AX=\beta$；记\textbf{增广矩阵}
\[
B=(A\mid\beta)=\left(\begin{array}{cccc|c}
a_{11}&a_{12}&\cdots&a_{1n}&b_1\\
a_{21}&a_{22}&\cdots&a_{2n}&b_2\\
\vdots&\vdots&&\vdots&\vdots\\
a_{m1}&a_{m2}&\cdots&a_{mn}&b_m
\end{array}\right).
\]
若把$A$按列分块为$A=(\alpha_1,\alpha_2,\cdots,\alpha_n)$，则方程组又可写成向量形式
\[
x_1\alpha_1+x_2\alpha_2+\cdots+x_n\alpha_n=\beta.
\]
于是“方程组有解”$\iff$“$\beta$可由$A$的列向量组线性表示”。若方程组有解，称它是\textbf{相容的}，否则称\textbf{不相容}。若$\beta\neq0$，称方程组为\textbf{非齐次}的；若$\beta=0$，即$AX=0$，称为\textbf{齐次}线性方程组（它总有零解，故总是相容的）。

\begin{theorem}[有解判定]
$n$元线性方程组$AX=\beta$有解的充要条件是
\[
R(A)=R(A\mid\beta).
\]
\end{theorem}

\begin{theorem}[解的个数]
设$R(A)=R(A\mid\beta)=r$，$n$为未知量个数，则
\begin{enumerate}
\item 若$r=n$，方程组有\textbf{唯一解}；
\item 若$r<n$，方程组有\textbf{无穷多解}，其中有$n-r$个自由未知量；
\item 若$R(A)<R(A\mid\beta)$，方程组\textbf{无解}。
\end{enumerate}
\end{theorem}

\keyformula{R(A)=R(A\mid\beta)=r:\quad r=n\Rightarrow\text{唯一解},\qquad r<n\Rightarrow\text{无穷多解};\qquad
R(A)<R(A\mid\beta)\Rightarrow\text{无解}}

判定定理的直观解释：把增广矩阵化为行阶梯形后，若出现形如$(0\ \cdots\ 0\mid d)$（$d\neq0$）的行，即“$0=d$”这样的矛盾方程，则无解——这正是$R(A)<R(A\mid\beta)$；若不出现矛盾行，则剩下的等式都能解，自由未知量个数就是$n-r$。

\subsection{齐次线性方程组}

对$AX=0$（$A$为$m\times n$矩阵），因$R(A)\le R(A\mid 0)=R(A)$恒成立，故\textbf{齐次方程组总是相容}（至少有零解）。

\begin{theorem}
设$A$为$m\times n$矩阵，则
\begin{enumerate}
\item $AX=0$只有零解$\iff R(A)=n$（$A$列满秩）；
\item $AX=0$有非零解$\iff R(A)<n$；
\item 特别地，当$m<n$（方程个数少于未知量个数）时，$AX=0$必有非零解；
\item 当$m=n$时，$AX=0$有非零解$\iff|A|=0$（与克拉默法则的推论一致）。
\end{enumerate}
\end{theorem}

当$R(A)=r<n$时，$AX=0$的全部解可表示为
\[
X=k_1\xi_1+k_2\xi_2+\cdots+k_{n-r}\xi_{n-r}\quad(k_1,\cdots,k_{n-r}\text{为任意常数}),
\]
其中$\xi_1,\cdots,\xi_{n-r}$是$AX=0$的$n-r$个线性无关的解向量（\textbf{基础解系}）。求法即下面的“自由未知量法”；基础解系与解空间的完整理论见第4章§5。

\subsection{非齐次线性方程组}

设$AX=\beta$（$\beta\neq0$），称$AX=0$为它的\textbf{导出组}（对应的齐次方程组）。

\begin{theorem}
设$\eta_1,\eta_2$都是$AX=\beta$的解，$\xi$是$AX=0$的解，则
\begin{enumerate}
\item $\eta_1-\eta_2$是$AX=0$的解；
\item $\eta_1+\xi$仍是$AX=\beta$的解。
\end{enumerate}
\end{theorem}

\begin{corollary}[通解结构]
若$AX=\beta$有解，$R(A)=r<n$，$\eta^{*}$是$AX=\beta$的一个\textbf{特解}，$\xi_1,\cdots,\xi_{n-r}$是导出组$AX=0$的基础解系，则$AX=\beta$的\textbf{通解}为
\[
X=\eta^{*}+k_1\xi_1+k_2\xi_2+\cdots+k_{n-r}\xi_{n-r}\quad(k_1,\cdots,k_{n-r}\text{为任意常数}).
\]
即“非齐次通解＝一个特解＋导出组通解”。
\end{corollary}

\subsection{用初等变换求通解（步骤）}

\begin{enumerate}
\item 写出增广矩阵$B=(A\mid\beta)$，只用\textbf{初等行变换}化为行阶梯形，判断$R(A)$与$R(A\mid\beta)$：若$R(A)<R(A\mid\beta)$，无解，停止；
\item 继续化为\textbf{行最简形}（首非零元为$1$，所在列其余元为$0$）；
\item 由行最简形写出同解的方程组：首非零元对应的未知量用其余未知量表示；
\item 把\textbf{自由未知量}（非首非零元对应的未知量）依次取$(1,0,\cdots,0),(0,1,\cdots,0),\cdots$，并令常数项为$0$，得导出组的基础解系$\xi_1,\cdots,\xi_{n-r}$；再令自由未知量全取$0$，得特解$\eta^*$；
\item 写出通解$X=\eta^*+k_1\xi_1+\cdots+k_{n-r}\xi_{n-r}$。
\end{enumerate}

\corepoint{
（1）解方程组\textbf{只能做初等行变换}，绝不能做列变换（列变换相当于把未知量改名，通解会对不上）。
（2）判断有无解必须用$R(A)$与$R(A\mid\beta)$，只看$R(A)$会误判。
（3）$R(A)=R(A\mid\beta)<n$时“有无穷多解”，答案必须写出\textbf{通解}含$n-r$个任意常数，只写一个特解算错。
（4）齐次方程组$AX=0$的基础解系必须取$n-R(A)$个\textbf{线性无关}的解；若自由未知量取值组选取不当，可能取到线性相关的向量组。
（5）含参数的方程组，先化阶梯形，再对参数使主元为零的取值逐个讨论（此时秩可能下降，解的个数随之改变）。
}

\subsection{与矩阵方程、矩阵秩的联系}

\begin{enumerate}
\item $AX=B$（$B$为$m\times p$矩阵）可视为$p$个“系数矩阵相同、常数项不同”的方程组同时求解，用$(A\mid B)\xrightarrow{\text{行变换}}$行最简形处理；若$A$可逆，则$X=A^{-1}B$。
\item $AX=0$有非零解$\iff R(A)<n$，这条判据在第4章用于判定向量组线性相关，在第5章用于求特征向量（$(\lambda E-A)X=0$有非零解$\iff|\lambda E-A|=0$）。
\item $R(AB)\le\min\{R(A),R(B)\}$与$AB=O\Rightarrow R(A)+R(B)\le n$可由“$B$的每一列都是$AX=0$的解”直接看出。
\end{enumerate}
